% ============================================================
%  Section: RAG Module — TounsiLM
%  Include in the main document with:
%    % ============================================================
%  Section: RAG Module — TounsiLM
%  Include in the main document with:
%    % ============================================================
%  Section: RAG Module — TounsiLM
%  Include in the main document with:
%    % ============================================================
%  Section: RAG Module — TounsiLM
%  Include in the main document with:
%    \input{rapport_rag_section}
%  Required packages in the preamble:
%    \usepackage[utf8]{inputenc}
%    \usepackage[T1]{fontenc}
%    \usepackage[english]{babel}
%    \usepackage{graphicx}
%    \usepackage{booktabs}
%    \usepackage{tabularx}
%    \usepackage{listings}
%    \usepackage{xcolor}
%    \usepackage{tikz}
%    \usetikzlibrary{shapes.geometric, arrows.meta, positioning, fit}
%    \usepackage{amsmath}
%    \usepackage{hyperref}
% ============================================================

\section{Retrieval-Augmented Generation (RAG) Module}
\label{sec:rag}

\subsection{Introduction and Motivation}

Large language models (LLMs) suffer from a fundamental limitation: their knowledge is
frozen at training time and does not cover niche domains or low-resource languages.
Tunisian Arabic (\textit{Tounsi}) is precisely such a domain — poorly documented,
rich in loanwords from French, Classical Arabic and Berber, and subject to strong
regional and generational variation.

To address this limitation, we integrated a \textbf{Retrieval-Augmented Generation (RAG)}
module into the \textbf{TounsiLM-8b} model. The core idea is straightforward: before
generating a response, the system dynamically \textit{retrieves} the most relevant
entries from a specialized knowledge base and injects them into the LLM's context.
This grounds the generation in verified facts without requiring any retraining of the model.

% -------------------------------------------------------
\subsection{Overall Architecture}
\label{subsec:architecture}

The RAG pipeline is composed of six main stages, illustrated in
Figure~\ref{fig:rag_pipeline}.

\begin{figure}[htbp]
\centering
\begin{tikzpicture}[
    box/.style={rectangle, rounded corners=4pt, draw=black!70, fill=blue!8,
                minimum width=3.0cm, minimum height=0.85cm, align=center, font=\small},
    newbox/.style={rectangle, rounded corners=4pt, draw=black!70, fill=purple!10,
                minimum width=3.0cm, minimum height=0.85cm, align=center, font=\small},
    store/.style={cylinder, shape border rotate=90, draw=black!70, fill=orange!12,
                  minimum width=2.2cm, minimum height=1.3cm, align=center, font=\small},
    arrow/.style={-Stealth, thick},
    dasharrow/.style={-Stealth, thick, dashed},
    lbl/.style={font=\footnotesize\itshape, text=gray}
]

%--- Top row: inference flow (left -> right) ---%
\node[box]    (query)    at ( 0.0,  0.0) {User query\\{\small Arabic/Arabizi/French}};
\node[newbox] (rewriter) at ( 4.0,  0.0) {Query Rewriter\\+ Type Router};
\node[newbox] (hybrid)   at ( 8.0,  0.0) {Hybrid Retrieval\\BM25 + Semantic};
\node[store]  (chroma)   at (11.5,  0.0) {ChromaDB\\(vectors)};

%--- Bottom row: inference flow (right -> left) ---%
\node[box]    (merge)    at (11.5, -3.2) {Multi-query merge\\(best score)};
\node[box]    (context)  at ( 8.0, -3.2) {Context\\formatting};
\node[box]    (llm)      at ( 4.0, -3.2) {TounsiLM-8b\\(generation)};
\node[box]    (answer)   at ( 0.0, -3.2) {Answer +\\confidence level};

%--- Indexing column above ChromaDB (dashed) ---%
\node[box, fill=green!8] (indexer) at (11.5,  1.8) {Indexing\\(\texttt{passage:})};
\node[box, fill=green!8] (kb)      at (11.5,  3.3) {Knowledge base\\(JSON --- 1{,}647 entries)};

%--- Inference arrows ---%
\draw[arrow] (query)         -- (rewriter);
\draw[arrow] (rewriter)      -- node[above, lbl]{variants} (hybrid);
\draw[arrow] (hybrid)        -- (chroma);
\draw[arrow] (chroma.south)  -- (merge.north);
\draw[arrow] (merge)         -- (context);
\draw[arrow] (context)       -- (llm);
\draw[arrow] (llm)           -- (answer);
% Original question fed directly to LLM (left-side bypass)
\draw[arrow] (query.west) -- ++(-0.45,0) |- (llm.west);

%--- Indexing arrows ---%
\draw[dasharrow] (kb)      -- (indexer);
\draw[dasharrow] (indexer) -- node[right, lbl]{index} (chroma.north);

\end{tikzpicture}
\caption{RAG pipeline architecture. Purple boxes are newly introduced components.
Dashed path: one-time indexing phase. Solid path: inference flow.
The left-side arrow shows the original query fed directly to TounsiLM-8b alongside the retrieved context.}
\label{fig:rag_pipeline}
\end{figure}

% -------------------------------------------------------
\subsection{Knowledge Base}
\label{subsec:kb}

The knowledge base consists of JSON files organized by semantic category.
Every entry is validated against a Pydantic schema before indexing, ensuring
data integrity throughout the pipeline.

\subsubsection{Entry Types}

Table~\ref{tab:types} summarizes the eleven supported entry types, their associated
Pydantic schema, and the type-specific fields they introduce beyond the base schema.

\begin{table}[htbp]
\centering
\renewcommand{\arraystretch}{1.25}
\begin{tabularx}{\textwidth}{@{} l l X @{}}
\toprule
\textbf{Type} & \textbf{Schema} & \textbf{Type-specific fields} \\
\midrule
\texttt{expression}    & \texttt{ExpressionEntry}   & \texttt{origin}, \texttt{severity}, \texttt{gender\_sensitive} \\
\texttt{number\_slang} & \texttt{NumberSlangEntry}  & \texttt{msa\_equivalent} \\
\texttt{proverb}       & \texttt{ProverbEntry}      & \texttt{literal\_meaning}, \texttt{real\_meaning}, \texttt{when\_used}, \texttt{msa\_equivalent} \\
\texttt{food}          & \texttt{FoodEntry}         & \texttt{description}, \texttt{regional\_variation}, \texttt{when\_eaten}, \texttt{similar\_to} \\
\texttt{ingredient}    & \texttt{FoodEntry}         & (same as \texttt{food}) \\
\texttt{ritual}        & \texttt{RitualEntry}       & \texttt{occasion}, \texttt{expected\_response}, \texttt{tone} \\
\texttt{code\_switch}  & \texttt{CodeSwitchEntry}   & \texttt{origin\_language}, \texttt{origin\_word}, \texttt{tunisian\_equivalent}, \texttt{domain} \\
\texttt{tv\_series}    & \texttt{MediaEntry}        & \texttt{channel\_platform}, \texttt{era}, \texttt{cultural\_significance}, \texttt{common\_references} \\
\texttt{movie}         & \texttt{MediaEntry}        & (same as \texttt{tv\_series}) \\
\texttt{film}          & \texttt{MediaEntry}        & (same as \texttt{tv\_series}) \\
\texttt{color}         & \texttt{ColorEntry}        & \texttt{color\_family}, \texttt{cultural\_significance} \\
\bottomrule
\end{tabularx}
\caption{Knowledge base entry types and their Pydantic schemas.}
\label{tab:types}
\end{table}

\subsubsection{Base Schema}

All types inherit from \texttt{BaseEntry}, which enforces the following common fields:
\texttt{id}, \texttt{type}, \texttt{term\_arabic}, \texttt{term\_arabizi}, \texttt{meaning},
\texttt{meaning\_fr} (optional), \texttt{example}, \texttt{usage\_context},
\texttt{region}, \texttt{register}, \texttt{generation}, \texttt{scripts},
\texttt{source}, and \texttt{last\_updated}.

\subsubsection{Data Volume}

The current version of the knowledge base contains \textbf{1,647 valid entries}
across 9 JSON files, covering idiomatic expressions, proverbs, gastronomy, social
rituals, French--Arabic code-switching, media culture, and colors in the Tunisian dialect.

% -------------------------------------------------------
\subsection{Indexing Pipeline}
\label{subsec:indexing}

Indexing is performed in two steps: embedding text construction and insertion into
the vector store.

\subsubsection{Embedding Text Construction}

The \texttt{build\_embed\_text} module constructs a text for each entry by concatenating
its most informative fields. The selected fields vary by type: for a proverb,
\texttt{literal\_meaning} and \texttt{real\_meaning} are appended; for a
\texttt{code\_switch} entry, \texttt{origin\_language} and \texttt{tunisian\_equivalent}
are added.

Following the recommendations of the \texttt{multilingual-e5} model, every indexed
text is prefixed with \texttt{"passage: "}:

\begin{equation}
    t_{\text{index}} = \texttt{"passage: "} \;\|\; \bigoplus_{i} \text{field}_i
    \label{eq:passage_prefix}
\end{equation}

\noindent where $\|$ denotes concatenation and $\bigoplus$ the union of relevant fields.
This asymmetric prefix (\texttt{"passage:"} at indexing time, \texttt{"query:"} at
retrieval time) is essential for correctly exploiting the model's embedding space.

\subsubsection{Vector Store}

Vectors are persisted in \textbf{ChromaDB} using cosine similarity as the distance
metric. Indexing is incremental: identifiers already present in the collection are
skipped during updates. A deduplication mechanism using a \texttt{seen\_ids} set
prevents ID collisions across distinct JSON files.

% -------------------------------------------------------
\subsection{Query Rewriting and Routing}
\label{subsec:rewriting}

Arabizi --- the Latin-script transcription of Arabic widely used in online communication ---
has no standardized spelling. The same word can appear as \textit{barcha},
\textit{barca}, or \textit{bar5a}; \textit{7amra} and \textit{hamra} denote the same
color. A single query formulation therefore risks missing relevant entries that were
indexed under a different spelling.

\subsubsection{Query Rewriting}

The \texttt{QueryRewriter} module addresses this by generating up to three variants
of each query before retrieval:

\begin{enumerate}
    \item \textbf{Original query} --- preserved as-is.
    \item \textbf{Arabizi-normalized variant} --- common digit-to-letter substitutions
    are applied (Table~\ref{tab:arabizi}), converting e.g.\ \textit{7amra} to
    \textit{hamra} and \textit{5adra} to \textit{khadra}.
    \item \textbf{Cleaned variant} --- punctuation is stripped to handle
    question marks, commas, and Arabic \textit{tatweel}.
\end{enumerate}

\begin{table}[htbp]
\centering
\begin{tabular}{@{} c c l @{}}
\toprule
\textbf{Digit} & \textbf{Replacement} & \textbf{Arabic phoneme} \\
\midrule
3 & a        & ayn \\
7 & h        & ha \\
5 & kh       & kha \\
9 & q        & qaf \\
8 & gh       & ghain \\
2 & (dropped) & hamza \\
\bottomrule
\end{tabular}
\caption{Arabizi digit-to-letter normalization rules applied by the query rewriter.}
\label{tab:arabizi}
\end{table}

Each variant is retrieved independently. Results are then merged by keeping the
highest score for each document across all variants:

\begin{equation}
    s^*(d) = \max_{v \in \mathcal{V}} s(d, v)
    \label{eq:multiquery_merge}
\end{equation}

\noindent where $\mathcal{V}$ is the set of query variants and $s(d,v)$ is the
RRF-fused score of document $d$ for variant $v$.

\subsubsection{Query Routing}

The same \texttt{QueryRewriter} module inspects the query for type-specific keywords
and automatically sets the \texttt{entry\_type} filter before retrieval. For example,
a query containing \textit{harissa}, \textit{food}, or the Arabic term for cooking
routes to \texttt{entry\_type="food"}, restricting the search space to the relevant
category. This improves both precision and speed without any user intervention.

% -------------------------------------------------------
\subsection{Hybrid Retrieval}
\label{subsec:retrieval}

Retrieval combines two complementary methods: dense semantic search and sparse
keyword search (BM25). Each handles a different failure mode of the other.

\subsubsection{Semantic Retrieval}

At inference time, the query variant is prefixed with \texttt{"query: "} and encoded
by \texttt{multilingual-e5-base}:

\begin{equation}
    \mathbf{q} = \text{Encode}\!\left(\texttt{"query: "} \| q\right)
    \label{eq:query_encode}
\end{equation}

ChromaDB computes cosine similarity between $\mathbf{q}$ and all indexed passage
vectors, returning $3k$ candidates to leave room for RRF re-ordering. The relevance
score is:

\begin{equation}
    s_{\text{sem}}(d) = 1 - d_{\cos}(\mathbf{q},\, \mathbf{p}_d)
    \label{eq:sem_score}
\end{equation}

\subsubsection{BM25 Retrieval}

A \textbf{BM25Okapi} index is built at startup over all 1,647 documents (stripped of
the \texttt{"passage: "} prefix). At query time, BM25 scores are normalized to $[0,1]$:

\begin{equation}
    s_{\text{BM25}}(d) = \frac{\text{BM25}(q, d)}{\max_{d'} \text{BM25}(q, d')}
    \label{eq:bm25_score}
\end{equation}

BM25 excels at exact term matches --- when a user types a precise Arabizi word or a
French loanword that dense embeddings might misplace in the embedding space.

\subsubsection{Reciprocal Rank Fusion}

The two ranked lists are merged using \textbf{Reciprocal Rank Fusion (RRF)} with
per-source weights:

\begin{equation}
    \text{RRF}(d) = \frac{w_{\text{sem}}}{k_0 + r_{\text{sem}}(d)}
                  + \frac{w_{\text{BM25}}}{k_0 + r_{\text{BM25}}(d)}
    \label{eq:rrf}
\end{equation}

\noindent where $r_i(d)$ is the rank of document $d$ in ranking $i$ (1-indexed),
$k_0 = 60$ is the RRF smoothing constant, $w_{\text{sem}} = 0.6$ and
$w_{\text{BM25}} = 0.4$. The final score is normalized to $[0,1]$ and the top-$k$
documents are returned.

% -------------------------------------------------------
\subsection{Confidence Scoring}
\label{subsec:confidence}

After retrieval, an overall confidence score is computed from the mean of the
RRF-normalized scores:

\begin{equation}
    \bar{s} = \frac{1}{k} \sum_{i=1}^{k} s_i
    \label{eq:mean_score}
\end{equation}

The confidence level is then assigned according to three tiers:

\begin{table}[htbp]
\centering
\begin{tabular}{@{} c l l @{}}
\toprule
\textbf{Indicator} & \textbf{Level} & \textbf{Condition} \\
\midrule
\textcolor{green!60!black}{$\bullet$} & High   & $\bar{s} \geq 0.75$ \\
\textcolor{orange}{$\bullet$}         & Medium & $0.50 \leq \bar{s} < 0.75$ \\
\textcolor{red}{$\bullet$}            & Low    & $\bar{s} < 0.50$ \\
$\bullet$                             & None   & no results returned \\
\bottomrule
\end{tabular}
\caption{Confidence levels based on the mean RRF score $\bar{s}$.}
\label{tab:confidence}
\end{table}

% -------------------------------------------------------
\subsection{Context Formatting}
\label{subsec:context}

The retrieved entries are formatted into a structured text block injected into the
LLM prompt. Each line follows the format:

\begin{quote}
\texttt{[type | score: \textit{s}] arabic\_term (arabizi) : meaning | \textnormal{مثال:} example | \textnormal{الاستخدام:} context}
\end{quote}

To avoid exceeding the model's context window (4,096 tokens), truncation is performed
at the token level using the TounsiLM-8b tokenizer, with a default limit of
\textbf{1,500 tokens} for the retrieved context.

% -------------------------------------------------------
\subsection{TounsiLM-8b Interface}
\label{subsec:llm}

The \textbf{TounsiLM-8b} model (\texttt{alabenayed/TounsiLM-8b}) is an 8-billion-parameter
LLM fine-tuned for understanding and generating Tunisian dialect using the
Llama-2 instruction format. The configurable generation parameters are summarized
in Table~\ref{tab:gen_params}.

\begin{table}[htbp]
\centering
\begin{tabular}{@{} l c l @{}}
\toprule
\textbf{Parameter}         & \textbf{Default} & \textbf{Role} \\
\midrule
\texttt{max\_new\_tokens}    & 512   & Maximum response length \\
\texttt{temperature}         & 0.7   & Creativity / diversity of generated text \\
\texttt{top\_p}              & 0.9   & Nucleus sampling \\
\texttt{top\_k}              & 50    & Top-$k$ sampling \\
\texttt{repetition\_penalty} & 1.1   & Repetition penalty \\
\bottomrule
\end{tabular}
\caption{TounsiLM-8b generation parameters.}
\label{tab:gen_params}
\end{table}

% -------------------------------------------------------
\subsection{Summary and Future Work}
\label{subsec:summary}

The developed RAG module brings several contributions to the TounsiLM system:

\begin{itemize}
    \item \textbf{Factual grounding}: responses are enriched with definitions,
    examples, and usage contexts drawn directly from the knowledge base, reducing
    hallucinations on dialectal vocabulary.

    \item \textbf{Arabizi robustness}: the query rewriter normalizes digit-based
    Arabizi spellings and generates multiple query variants, significantly improving
    recall for non-standard user inputs.

    \item \textbf{Hybrid precision}: combining BM25 and semantic search via RRF
    captures both exact keyword matches and semantic similarity, outperforming
    either method alone.

    \item \textbf{Automatic routing}: type-keyword detection routes queries to the
    appropriate KB category without user intervention, improving precision on
    domain-specific questions.

    \item \textbf{Extensibility}: adding a new data type requires only creating a
    Pydantic schema and a single line in the \texttt{SCHEMA\_MAP} dispatch table.

    \item \textbf{Transparency}: confidence level and sources are systematically
    exposed to the user, facilitating qualitative evaluation of responses.
\end{itemize}

Among the planned improvements are a \textbf{re-ranking} step using a cross-encoder
to further refine the top-$k$ results before generation, and a
\textbf{history-aware} conversational wrapper to maintain context across multiple
turns in interactive mode.

%  Required packages in the preamble:
%    \usepackage[utf8]{inputenc}
%    \usepackage[T1]{fontenc}
%    \usepackage[english]{babel}
%    \usepackage{graphicx}
%    \usepackage{booktabs}
%    \usepackage{tabularx}
%    \usepackage{listings}
%    \usepackage{xcolor}
%    \usepackage{tikz}
%    \usetikzlibrary{shapes.geometric, arrows.meta, positioning, fit}
%    \usepackage{amsmath}
%    \usepackage{hyperref}
% ============================================================

\section{Retrieval-Augmented Generation (RAG) Module}
\label{sec:rag}

\subsection{Introduction and Motivation}

Large language models (LLMs) suffer from a fundamental limitation: their knowledge is
frozen at training time and does not cover niche domains or low-resource languages.
Tunisian Arabic (\textit{Tounsi}) is precisely such a domain — poorly documented,
rich in loanwords from French, Classical Arabic and Berber, and subject to strong
regional and generational variation.

To address this limitation, we integrated a \textbf{Retrieval-Augmented Generation (RAG)}
module into the \textbf{TounsiLM-8b} model. The core idea is straightforward: before
generating a response, the system dynamically \textit{retrieves} the most relevant
entries from a specialized knowledge base and injects them into the LLM's context.
This grounds the generation in verified facts without requiring any retraining of the model.

% -------------------------------------------------------
\subsection{Overall Architecture}
\label{subsec:architecture}

The RAG pipeline is composed of six main stages, illustrated in
Figure~\ref{fig:rag_pipeline}.

\begin{figure}[htbp]
\centering
\begin{tikzpicture}[
    box/.style={rectangle, rounded corners=4pt, draw=black!70, fill=blue!8,
                minimum width=3.0cm, minimum height=0.85cm, align=center, font=\small},
    newbox/.style={rectangle, rounded corners=4pt, draw=black!70, fill=purple!10,
                minimum width=3.0cm, minimum height=0.85cm, align=center, font=\small},
    store/.style={cylinder, shape border rotate=90, draw=black!70, fill=orange!12,
                  minimum width=2.2cm, minimum height=1.3cm, align=center, font=\small},
    arrow/.style={-Stealth, thick},
    dasharrow/.style={-Stealth, thick, dashed},
    lbl/.style={font=\footnotesize\itshape, text=gray}
]

%--- Top row: inference flow (left -> right) ---%
\node[box]    (query)    at ( 0.0,  0.0) {User query\\{\small Arabic/Arabizi/French}};
\node[newbox] (rewriter) at ( 4.0,  0.0) {Query Rewriter\\+ Type Router};
\node[newbox] (hybrid)   at ( 8.0,  0.0) {Hybrid Retrieval\\BM25 + Semantic};
\node[store]  (chroma)   at (11.5,  0.0) {ChromaDB\\(vectors)};

%--- Bottom row: inference flow (right -> left) ---%
\node[box]    (merge)    at (11.5, -3.2) {Multi-query merge\\(best score)};
\node[box]    (context)  at ( 8.0, -3.2) {Context\\formatting};
\node[box]    (llm)      at ( 4.0, -3.2) {TounsiLM-8b\\(generation)};
\node[box]    (answer)   at ( 0.0, -3.2) {Answer +\\confidence level};

%--- Indexing column above ChromaDB (dashed) ---%
\node[box, fill=green!8] (indexer) at (11.5,  1.8) {Indexing\\(\texttt{passage:})};
\node[box, fill=green!8] (kb)      at (11.5,  3.3) {Knowledge base\\(JSON --- 1{,}647 entries)};

%--- Inference arrows ---%
\draw[arrow] (query)         -- (rewriter);
\draw[arrow] (rewriter)      -- node[above, lbl]{variants} (hybrid);
\draw[arrow] (hybrid)        -- (chroma);
\draw[arrow] (chroma.south)  -- (merge.north);
\draw[arrow] (merge)         -- (context);
\draw[arrow] (context)       -- (llm);
\draw[arrow] (llm)           -- (answer);
% Original question fed directly to LLM (left-side bypass)
\draw[arrow] (query.west) -- ++(-0.45,0) |- (llm.west);

%--- Indexing arrows ---%
\draw[dasharrow] (kb)      -- (indexer);
\draw[dasharrow] (indexer) -- node[right, lbl]{index} (chroma.north);

\end{tikzpicture}
\caption{RAG pipeline architecture. Purple boxes are newly introduced components.
Dashed path: one-time indexing phase. Solid path: inference flow.
The left-side arrow shows the original query fed directly to TounsiLM-8b alongside the retrieved context.}
\label{fig:rag_pipeline}
\end{figure}

% -------------------------------------------------------
\subsection{Knowledge Base}
\label{subsec:kb}

The knowledge base consists of JSON files organized by semantic category.
Every entry is validated against a Pydantic schema before indexing, ensuring
data integrity throughout the pipeline.

\subsubsection{Entry Types}

Table~\ref{tab:types} summarizes the eleven supported entry types, their associated
Pydantic schema, and the type-specific fields they introduce beyond the base schema.

\begin{table}[htbp]
\centering
\renewcommand{\arraystretch}{1.25}
\begin{tabularx}{\textwidth}{@{} l l X @{}}
\toprule
\textbf{Type} & \textbf{Schema} & \textbf{Type-specific fields} \\
\midrule
\texttt{expression}    & \texttt{ExpressionEntry}   & \texttt{origin}, \texttt{severity}, \texttt{gender\_sensitive} \\
\texttt{number\_slang} & \texttt{NumberSlangEntry}  & \texttt{msa\_equivalent} \\
\texttt{proverb}       & \texttt{ProverbEntry}      & \texttt{literal\_meaning}, \texttt{real\_meaning}, \texttt{when\_used}, \texttt{msa\_equivalent} \\
\texttt{food}          & \texttt{FoodEntry}         & \texttt{description}, \texttt{regional\_variation}, \texttt{when\_eaten}, \texttt{similar\_to} \\
\texttt{ingredient}    & \texttt{FoodEntry}         & (same as \texttt{food}) \\
\texttt{ritual}        & \texttt{RitualEntry}       & \texttt{occasion}, \texttt{expected\_response}, \texttt{tone} \\
\texttt{code\_switch}  & \texttt{CodeSwitchEntry}   & \texttt{origin\_language}, \texttt{origin\_word}, \texttt{tunisian\_equivalent}, \texttt{domain} \\
\texttt{tv\_series}    & \texttt{MediaEntry}        & \texttt{channel\_platform}, \texttt{era}, \texttt{cultural\_significance}, \texttt{common\_references} \\
\texttt{movie}         & \texttt{MediaEntry}        & (same as \texttt{tv\_series}) \\
\texttt{film}          & \texttt{MediaEntry}        & (same as \texttt{tv\_series}) \\
\texttt{color}         & \texttt{ColorEntry}        & \texttt{color\_family}, \texttt{cultural\_significance} \\
\bottomrule
\end{tabularx}
\caption{Knowledge base entry types and their Pydantic schemas.}
\label{tab:types}
\end{table}

\subsubsection{Base Schema}

All types inherit from \texttt{BaseEntry}, which enforces the following common fields:
\texttt{id}, \texttt{type}, \texttt{term\_arabic}, \texttt{term\_arabizi}, \texttt{meaning},
\texttt{meaning\_fr} (optional), \texttt{example}, \texttt{usage\_context},
\texttt{region}, \texttt{register}, \texttt{generation}, \texttt{scripts},
\texttt{source}, and \texttt{last\_updated}.

\subsubsection{Data Volume}

The current version of the knowledge base contains \textbf{1,647 valid entries}
across 9 JSON files, covering idiomatic expressions, proverbs, gastronomy, social
rituals, French--Arabic code-switching, media culture, and colors in the Tunisian dialect.

% -------------------------------------------------------
\subsection{Indexing Pipeline}
\label{subsec:indexing}

Indexing is performed in two steps: embedding text construction and insertion into
the vector store.

\subsubsection{Embedding Text Construction}

The \texttt{build\_embed\_text} module constructs a text for each entry by concatenating
its most informative fields. The selected fields vary by type: for a proverb,
\texttt{literal\_meaning} and \texttt{real\_meaning} are appended; for a
\texttt{code\_switch} entry, \texttt{origin\_language} and \texttt{tunisian\_equivalent}
are added.

Following the recommendations of the \texttt{multilingual-e5} model, every indexed
text is prefixed with \texttt{"passage: "}:

\begin{equation}
    t_{\text{index}} = \texttt{"passage: "} \;\|\; \bigoplus_{i} \text{field}_i
    \label{eq:passage_prefix}
\end{equation}

\noindent where $\|$ denotes concatenation and $\bigoplus$ the union of relevant fields.
This asymmetric prefix (\texttt{"passage:"} at indexing time, \texttt{"query:"} at
retrieval time) is essential for correctly exploiting the model's embedding space.

\subsubsection{Vector Store}

Vectors are persisted in \textbf{ChromaDB} using cosine similarity as the distance
metric. Indexing is incremental: identifiers already present in the collection are
skipped during updates. A deduplication mechanism using a \texttt{seen\_ids} set
prevents ID collisions across distinct JSON files.

% -------------------------------------------------------
\subsection{Query Rewriting and Routing}
\label{subsec:rewriting}

Arabizi --- the Latin-script transcription of Arabic widely used in online communication ---
has no standardized spelling. The same word can appear as \textit{barcha},
\textit{barca}, or \textit{bar5a}; \textit{7amra} and \textit{hamra} denote the same
color. A single query formulation therefore risks missing relevant entries that were
indexed under a different spelling.

\subsubsection{Query Rewriting}

The \texttt{QueryRewriter} module addresses this by generating up to three variants
of each query before retrieval:

\begin{enumerate}
    \item \textbf{Original query} --- preserved as-is.
    \item \textbf{Arabizi-normalized variant} --- common digit-to-letter substitutions
    are applied (Table~\ref{tab:arabizi}), converting e.g.\ \textit{7amra} to
    \textit{hamra} and \textit{5adra} to \textit{khadra}.
    \item \textbf{Cleaned variant} --- punctuation is stripped to handle
    question marks, commas, and Arabic \textit{tatweel}.
\end{enumerate}

\begin{table}[htbp]
\centering
\begin{tabular}{@{} c c l @{}}
\toprule
\textbf{Digit} & \textbf{Replacement} & \textbf{Arabic phoneme} \\
\midrule
3 & a        & ayn \\
7 & h        & ha \\
5 & kh       & kha \\
9 & q        & qaf \\
8 & gh       & ghain \\
2 & (dropped) & hamza \\
\bottomrule
\end{tabular}
\caption{Arabizi digit-to-letter normalization rules applied by the query rewriter.}
\label{tab:arabizi}
\end{table}

Each variant is retrieved independently. Results are then merged by keeping the
highest score for each document across all variants:

\begin{equation}
    s^*(d) = \max_{v \in \mathcal{V}} s(d, v)
    \label{eq:multiquery_merge}
\end{equation}

\noindent where $\mathcal{V}$ is the set of query variants and $s(d,v)$ is the
RRF-fused score of document $d$ for variant $v$.

\subsubsection{Query Routing}

The same \texttt{QueryRewriter} module inspects the query for type-specific keywords
and automatically sets the \texttt{entry\_type} filter before retrieval. For example,
a query containing \textit{harissa}, \textit{food}, or the Arabic term for cooking
routes to \texttt{entry\_type="food"}, restricting the search space to the relevant
category. This improves both precision and speed without any user intervention.

% -------------------------------------------------------
\subsection{Hybrid Retrieval}
\label{subsec:retrieval}

Retrieval combines two complementary methods: dense semantic search and sparse
keyword search (BM25). Each handles a different failure mode of the other.

\subsubsection{Semantic Retrieval}

At inference time, the query variant is prefixed with \texttt{"query: "} and encoded
by \texttt{multilingual-e5-base}:

\begin{equation}
    \mathbf{q} = \text{Encode}\!\left(\texttt{"query: "} \| q\right)
    \label{eq:query_encode}
\end{equation}

ChromaDB computes cosine similarity between $\mathbf{q}$ and all indexed passage
vectors, returning $3k$ candidates to leave room for RRF re-ordering. The relevance
score is:

\begin{equation}
    s_{\text{sem}}(d) = 1 - d_{\cos}(\mathbf{q},\, \mathbf{p}_d)
    \label{eq:sem_score}
\end{equation}

\subsubsection{BM25 Retrieval}

A \textbf{BM25Okapi} index is built at startup over all 1,647 documents (stripped of
the \texttt{"passage: "} prefix). At query time, BM25 scores are normalized to $[0,1]$:

\begin{equation}
    s_{\text{BM25}}(d) = \frac{\text{BM25}(q, d)}{\max_{d'} \text{BM25}(q, d')}
    \label{eq:bm25_score}
\end{equation}

BM25 excels at exact term matches --- when a user types a precise Arabizi word or a
French loanword that dense embeddings might misplace in the embedding space.

\subsubsection{Reciprocal Rank Fusion}

The two ranked lists are merged using \textbf{Reciprocal Rank Fusion (RRF)} with
per-source weights:

\begin{equation}
    \text{RRF}(d) = \frac{w_{\text{sem}}}{k_0 + r_{\text{sem}}(d)}
                  + \frac{w_{\text{BM25}}}{k_0 + r_{\text{BM25}}(d)}
    \label{eq:rrf}
\end{equation}

\noindent where $r_i(d)$ is the rank of document $d$ in ranking $i$ (1-indexed),
$k_0 = 60$ is the RRF smoothing constant, $w_{\text{sem}} = 0.6$ and
$w_{\text{BM25}} = 0.4$. The final score is normalized to $[0,1]$ and the top-$k$
documents are returned.

% -------------------------------------------------------
\subsection{Confidence Scoring}
\label{subsec:confidence}

After retrieval, an overall confidence score is computed from the mean of the
RRF-normalized scores:

\begin{equation}
    \bar{s} = \frac{1}{k} \sum_{i=1}^{k} s_i
    \label{eq:mean_score}
\end{equation}

The confidence level is then assigned according to three tiers:

\begin{table}[htbp]
\centering
\begin{tabular}{@{} c l l @{}}
\toprule
\textbf{Indicator} & \textbf{Level} & \textbf{Condition} \\
\midrule
\textcolor{green!60!black}{$\bullet$} & High   & $\bar{s} \geq 0.75$ \\
\textcolor{orange}{$\bullet$}         & Medium & $0.50 \leq \bar{s} < 0.75$ \\
\textcolor{red}{$\bullet$}            & Low    & $\bar{s} < 0.50$ \\
$\bullet$                             & None   & no results returned \\
\bottomrule
\end{tabular}
\caption{Confidence levels based on the mean RRF score $\bar{s}$.}
\label{tab:confidence}
\end{table}

% -------------------------------------------------------
\subsection{Context Formatting}
\label{subsec:context}

The retrieved entries are formatted into a structured text block injected into the
LLM prompt. Each line follows the format:

\begin{quote}
\texttt{[type | score: \textit{s}] arabic\_term (arabizi) : meaning | \textnormal{مثال:} example | \textnormal{الاستخدام:} context}
\end{quote}

To avoid exceeding the model's context window (4,096 tokens), truncation is performed
at the token level using the TounsiLM-8b tokenizer, with a default limit of
\textbf{1,500 tokens} for the retrieved context.

% -------------------------------------------------------
\subsection{TounsiLM-8b Interface}
\label{subsec:llm}

The \textbf{TounsiLM-8b} model (\texttt{alabenayed/TounsiLM-8b}) is an 8-billion-parameter
LLM fine-tuned for understanding and generating Tunisian dialect using the
Llama-2 instruction format. The configurable generation parameters are summarized
in Table~\ref{tab:gen_params}.

\begin{table}[htbp]
\centering
\begin{tabular}{@{} l c l @{}}
\toprule
\textbf{Parameter}         & \textbf{Default} & \textbf{Role} \\
\midrule
\texttt{max\_new\_tokens}    & 512   & Maximum response length \\
\texttt{temperature}         & 0.7   & Creativity / diversity of generated text \\
\texttt{top\_p}              & 0.9   & Nucleus sampling \\
\texttt{top\_k}              & 50    & Top-$k$ sampling \\
\texttt{repetition\_penalty} & 1.1   & Repetition penalty \\
\bottomrule
\end{tabular}
\caption{TounsiLM-8b generation parameters.}
\label{tab:gen_params}
\end{table}

% -------------------------------------------------------
\subsection{Summary and Future Work}
\label{subsec:summary}

The developed RAG module brings several contributions to the TounsiLM system:

\begin{itemize}
    \item \textbf{Factual grounding}: responses are enriched with definitions,
    examples, and usage contexts drawn directly from the knowledge base, reducing
    hallucinations on dialectal vocabulary.

    \item \textbf{Arabizi robustness}: the query rewriter normalizes digit-based
    Arabizi spellings and generates multiple query variants, significantly improving
    recall for non-standard user inputs.

    \item \textbf{Hybrid precision}: combining BM25 and semantic search via RRF
    captures both exact keyword matches and semantic similarity, outperforming
    either method alone.

    \item \textbf{Automatic routing}: type-keyword detection routes queries to the
    appropriate KB category without user intervention, improving precision on
    domain-specific questions.

    \item \textbf{Extensibility}: adding a new data type requires only creating a
    Pydantic schema and a single line in the \texttt{SCHEMA\_MAP} dispatch table.

    \item \textbf{Transparency}: confidence level and sources are systematically
    exposed to the user, facilitating qualitative evaluation of responses.
\end{itemize}

Among the planned improvements are a \textbf{re-ranking} step using a cross-encoder
to further refine the top-$k$ results before generation, and a
\textbf{history-aware} conversational wrapper to maintain context across multiple
turns in interactive mode.

%  Required packages in the preamble:
%    \usepackage[utf8]{inputenc}
%    \usepackage[T1]{fontenc}
%    \usepackage[english]{babel}
%    \usepackage{graphicx}
%    \usepackage{booktabs}
%    \usepackage{tabularx}
%    \usepackage{listings}
%    \usepackage{xcolor}
%    \usepackage{tikz}
%    \usetikzlibrary{shapes.geometric, arrows.meta, positioning, fit}
%    \usepackage{amsmath}
%    \usepackage{hyperref}
% ============================================================

\section{Retrieval-Augmented Generation (RAG) Module}
\label{sec:rag}

\subsection{Introduction and Motivation}

Large language models (LLMs) suffer from a fundamental limitation: their knowledge is
frozen at training time and does not cover niche domains or low-resource languages.
Tunisian Arabic (\textit{Tounsi}) is precisely such a domain — poorly documented,
rich in loanwords from French, Classical Arabic and Berber, and subject to strong
regional and generational variation.

To address this limitation, we integrated a \textbf{Retrieval-Augmented Generation (RAG)}
module into the \textbf{TounsiLM-8b} model. The core idea is straightforward: before
generating a response, the system dynamically \textit{retrieves} the most relevant
entries from a specialized knowledge base and injects them into the LLM's context.
This grounds the generation in verified facts without requiring any retraining of the model.

% -------------------------------------------------------
\subsection{Overall Architecture}
\label{subsec:architecture}

The RAG pipeline is composed of six main stages, illustrated in
Figure~\ref{fig:rag_pipeline}.

\begin{figure}[htbp]
\centering
\begin{tikzpicture}[
    box/.style={rectangle, rounded corners=4pt, draw=black!70, fill=blue!8,
                minimum width=3.0cm, minimum height=0.85cm, align=center, font=\small},
    newbox/.style={rectangle, rounded corners=4pt, draw=black!70, fill=purple!10,
                minimum width=3.0cm, minimum height=0.85cm, align=center, font=\small},
    store/.style={cylinder, shape border rotate=90, draw=black!70, fill=orange!12,
                  minimum width=2.2cm, minimum height=1.3cm, align=center, font=\small},
    arrow/.style={-Stealth, thick},
    dasharrow/.style={-Stealth, thick, dashed},
    lbl/.style={font=\footnotesize\itshape, text=gray}
]

%--- Top row: inference flow (left -> right) ---%
\node[box]    (query)    at ( 0.0,  0.0) {User query\\{\small Arabic/Arabizi/French}};
\node[newbox] (rewriter) at ( 4.0,  0.0) {Query Rewriter\\+ Type Router};
\node[newbox] (hybrid)   at ( 8.0,  0.0) {Hybrid Retrieval\\BM25 + Semantic};
\node[store]  (chroma)   at (11.5,  0.0) {ChromaDB\\(vectors)};

%--- Bottom row: inference flow (right -> left) ---%
\node[box]    (merge)    at (11.5, -3.2) {Multi-query merge\\(best score)};
\node[box]    (context)  at ( 8.0, -3.2) {Context\\formatting};
\node[box]    (llm)      at ( 4.0, -3.2) {TounsiLM-8b\\(generation)};
\node[box]    (answer)   at ( 0.0, -3.2) {Answer +\\confidence level};

%--- Indexing column above ChromaDB (dashed) ---%
\node[box, fill=green!8] (indexer) at (11.5,  1.8) {Indexing\\(\texttt{passage:})};
\node[box, fill=green!8] (kb)      at (11.5,  3.3) {Knowledge base\\(JSON --- 1{,}647 entries)};

%--- Inference arrows ---%
\draw[arrow] (query)         -- (rewriter);
\draw[arrow] (rewriter)      -- node[above, lbl]{variants} (hybrid);
\draw[arrow] (hybrid)        -- (chroma);
\draw[arrow] (chroma.south)  -- (merge.north);
\draw[arrow] (merge)         -- (context);
\draw[arrow] (context)       -- (llm);
\draw[arrow] (llm)           -- (answer);
% Original question fed directly to LLM (left-side bypass)
\draw[arrow] (query.west) -- ++(-0.45,0) |- (llm.west);

%--- Indexing arrows ---%
\draw[dasharrow] (kb)      -- (indexer);
\draw[dasharrow] (indexer) -- node[right, lbl]{index} (chroma.north);

\end{tikzpicture}
\caption{RAG pipeline architecture. Purple boxes are newly introduced components.
Dashed path: one-time indexing phase. Solid path: inference flow.
The left-side arrow shows the original query fed directly to TounsiLM-8b alongside the retrieved context.}
\label{fig:rag_pipeline}
\end{figure}

% -------------------------------------------------------
\subsection{Knowledge Base}
\label{subsec:kb}

The knowledge base consists of JSON files organized by semantic category.
Every entry is validated against a Pydantic schema before indexing, ensuring
data integrity throughout the pipeline.

\subsubsection{Entry Types}

Table~\ref{tab:types} summarizes the eleven supported entry types, their associated
Pydantic schema, and the type-specific fields they introduce beyond the base schema.

\begin{table}[htbp]
\centering
\renewcommand{\arraystretch}{1.25}
\begin{tabularx}{\textwidth}{@{} l l X @{}}
\toprule
\textbf{Type} & \textbf{Schema} & \textbf{Type-specific fields} \\
\midrule
\texttt{expression}    & \texttt{ExpressionEntry}   & \texttt{origin}, \texttt{severity}, \texttt{gender\_sensitive} \\
\texttt{number\_slang} & \texttt{NumberSlangEntry}  & \texttt{msa\_equivalent} \\
\texttt{proverb}       & \texttt{ProverbEntry}      & \texttt{literal\_meaning}, \texttt{real\_meaning}, \texttt{when\_used}, \texttt{msa\_equivalent} \\
\texttt{food}          & \texttt{FoodEntry}         & \texttt{description}, \texttt{regional\_variation}, \texttt{when\_eaten}, \texttt{similar\_to} \\
\texttt{ingredient}    & \texttt{FoodEntry}         & (same as \texttt{food}) \\
\texttt{ritual}        & \texttt{RitualEntry}       & \texttt{occasion}, \texttt{expected\_response}, \texttt{tone} \\
\texttt{code\_switch}  & \texttt{CodeSwitchEntry}   & \texttt{origin\_language}, \texttt{origin\_word}, \texttt{tunisian\_equivalent}, \texttt{domain} \\
\texttt{tv\_series}    & \texttt{MediaEntry}        & \texttt{channel\_platform}, \texttt{era}, \texttt{cultural\_significance}, \texttt{common\_references} \\
\texttt{movie}         & \texttt{MediaEntry}        & (same as \texttt{tv\_series}) \\
\texttt{film}          & \texttt{MediaEntry}        & (same as \texttt{tv\_series}) \\
\texttt{color}         & \texttt{ColorEntry}        & \texttt{color\_family}, \texttt{cultural\_significance} \\
\bottomrule
\end{tabularx}
\caption{Knowledge base entry types and their Pydantic schemas.}
\label{tab:types}
\end{table}

\subsubsection{Base Schema}

All types inherit from \texttt{BaseEntry}, which enforces the following common fields:
\texttt{id}, \texttt{type}, \texttt{term\_arabic}, \texttt{term\_arabizi}, \texttt{meaning},
\texttt{meaning\_fr} (optional), \texttt{example}, \texttt{usage\_context},
\texttt{region}, \texttt{register}, \texttt{generation}, \texttt{scripts},
\texttt{source}, and \texttt{last\_updated}.

\subsubsection{Data Volume}

The current version of the knowledge base contains \textbf{1,647 valid entries}
across 9 JSON files, covering idiomatic expressions, proverbs, gastronomy, social
rituals, French--Arabic code-switching, media culture, and colors in the Tunisian dialect.

% -------------------------------------------------------
\subsection{Indexing Pipeline}
\label{subsec:indexing}

Indexing is performed in two steps: embedding text construction and insertion into
the vector store.

\subsubsection{Embedding Text Construction}

The \texttt{build\_embed\_text} module constructs a text for each entry by concatenating
its most informative fields. The selected fields vary by type: for a proverb,
\texttt{literal\_meaning} and \texttt{real\_meaning} are appended; for a
\texttt{code\_switch} entry, \texttt{origin\_language} and \texttt{tunisian\_equivalent}
are added.

Following the recommendations of the \texttt{multilingual-e5} model, every indexed
text is prefixed with \texttt{"passage: "}:

\begin{equation}
    t_{\text{index}} = \texttt{"passage: "} \;\|\; \bigoplus_{i} \text{field}_i
    \label{eq:passage_prefix}
\end{equation}

\noindent where $\|$ denotes concatenation and $\bigoplus$ the union of relevant fields.
This asymmetric prefix (\texttt{"passage:"} at indexing time, \texttt{"query:"} at
retrieval time) is essential for correctly exploiting the model's embedding space.

\subsubsection{Vector Store}

Vectors are persisted in \textbf{ChromaDB} using cosine similarity as the distance
metric. Indexing is incremental: identifiers already present in the collection are
skipped during updates. A deduplication mechanism using a \texttt{seen\_ids} set
prevents ID collisions across distinct JSON files.

% -------------------------------------------------------
\subsection{Query Rewriting and Routing}
\label{subsec:rewriting}

Arabizi --- the Latin-script transcription of Arabic widely used in online communication ---
has no standardized spelling. The same word can appear as \textit{barcha},
\textit{barca}, or \textit{bar5a}; \textit{7amra} and \textit{hamra} denote the same
color. A single query formulation therefore risks missing relevant entries that were
indexed under a different spelling.

\subsubsection{Query Rewriting}

The \texttt{QueryRewriter} module addresses this by generating up to three variants
of each query before retrieval:

\begin{enumerate}
    \item \textbf{Original query} --- preserved as-is.
    \item \textbf{Arabizi-normalized variant} --- common digit-to-letter substitutions
    are applied (Table~\ref{tab:arabizi}), converting e.g.\ \textit{7amra} to
    \textit{hamra} and \textit{5adra} to \textit{khadra}.
    \item \textbf{Cleaned variant} --- punctuation is stripped to handle
    question marks, commas, and Arabic \textit{tatweel}.
\end{enumerate}

\begin{table}[htbp]
\centering
\begin{tabular}{@{} c c l @{}}
\toprule
\textbf{Digit} & \textbf{Replacement} & \textbf{Arabic phoneme} \\
\midrule
3 & a        & ayn \\
7 & h        & ha \\
5 & kh       & kha \\
9 & q        & qaf \\
8 & gh       & ghain \\
2 & (dropped) & hamza \\
\bottomrule
\end{tabular}
\caption{Arabizi digit-to-letter normalization rules applied by the query rewriter.}
\label{tab:arabizi}
\end{table}

Each variant is retrieved independently. Results are then merged by keeping the
highest score for each document across all variants:

\begin{equation}
    s^*(d) = \max_{v \in \mathcal{V}} s(d, v)
    \label{eq:multiquery_merge}
\end{equation}

\noindent where $\mathcal{V}$ is the set of query variants and $s(d,v)$ is the
RRF-fused score of document $d$ for variant $v$.

\subsubsection{Query Routing}

The same \texttt{QueryRewriter} module inspects the query for type-specific keywords
and automatically sets the \texttt{entry\_type} filter before retrieval. For example,
a query containing \textit{harissa}, \textit{food}, or the Arabic term for cooking
routes to \texttt{entry\_type="food"}, restricting the search space to the relevant
category. This improves both precision and speed without any user intervention.

% -------------------------------------------------------
\subsection{Hybrid Retrieval}
\label{subsec:retrieval}

Retrieval combines two complementary methods: dense semantic search and sparse
keyword search (BM25). Each handles a different failure mode of the other.

\subsubsection{Semantic Retrieval}

At inference time, the query variant is prefixed with \texttt{"query: "} and encoded
by \texttt{multilingual-e5-base}:

\begin{equation}
    \mathbf{q} = \text{Encode}\!\left(\texttt{"query: "} \| q\right)
    \label{eq:query_encode}
\end{equation}

ChromaDB computes cosine similarity between $\mathbf{q}$ and all indexed passage
vectors, returning $3k$ candidates to leave room for RRF re-ordering. The relevance
score is:

\begin{equation}
    s_{\text{sem}}(d) = 1 - d_{\cos}(\mathbf{q},\, \mathbf{p}_d)
    \label{eq:sem_score}
\end{equation}

\subsubsection{BM25 Retrieval}

A \textbf{BM25Okapi} index is built at startup over all 1,647 documents (stripped of
the \texttt{"passage: "} prefix). At query time, BM25 scores are normalized to $[0,1]$:

\begin{equation}
    s_{\text{BM25}}(d) = \frac{\text{BM25}(q, d)}{\max_{d'} \text{BM25}(q, d')}
    \label{eq:bm25_score}
\end{equation}

BM25 excels at exact term matches --- when a user types a precise Arabizi word or a
French loanword that dense embeddings might misplace in the embedding space.

\subsubsection{Reciprocal Rank Fusion}

The two ranked lists are merged using \textbf{Reciprocal Rank Fusion (RRF)} with
per-source weights:

\begin{equation}
    \text{RRF}(d) = \frac{w_{\text{sem}}}{k_0 + r_{\text{sem}}(d)}
                  + \frac{w_{\text{BM25}}}{k_0 + r_{\text{BM25}}(d)}
    \label{eq:rrf}
\end{equation}

\noindent where $r_i(d)$ is the rank of document $d$ in ranking $i$ (1-indexed),
$k_0 = 60$ is the RRF smoothing constant, $w_{\text{sem}} = 0.6$ and
$w_{\text{BM25}} = 0.4$. The final score is normalized to $[0,1]$ and the top-$k$
documents are returned.

% -------------------------------------------------------
\subsection{Confidence Scoring}
\label{subsec:confidence}

After retrieval, an overall confidence score is computed from the mean of the
RRF-normalized scores:

\begin{equation}
    \bar{s} = \frac{1}{k} \sum_{i=1}^{k} s_i
    \label{eq:mean_score}
\end{equation}

The confidence level is then assigned according to three tiers:

\begin{table}[htbp]
\centering
\begin{tabular}{@{} c l l @{}}
\toprule
\textbf{Indicator} & \textbf{Level} & \textbf{Condition} \\
\midrule
\textcolor{green!60!black}{$\bullet$} & High   & $\bar{s} \geq 0.75$ \\
\textcolor{orange}{$\bullet$}         & Medium & $0.50 \leq \bar{s} < 0.75$ \\
\textcolor{red}{$\bullet$}            & Low    & $\bar{s} < 0.50$ \\
$\bullet$                             & None   & no results returned \\
\bottomrule
\end{tabular}
\caption{Confidence levels based on the mean RRF score $\bar{s}$.}
\label{tab:confidence}
\end{table}

% -------------------------------------------------------
\subsection{Context Formatting}
\label{subsec:context}

The retrieved entries are formatted into a structured text block injected into the
LLM prompt. Each line follows the format:

\begin{quote}
\texttt{[type | score: \textit{s}] arabic\_term (arabizi) : meaning | \textnormal{مثال:} example | \textnormal{الاستخدام:} context}
\end{quote}

To avoid exceeding the model's context window (4,096 tokens), truncation is performed
at the token level using the TounsiLM-8b tokenizer, with a default limit of
\textbf{1,500 tokens} for the retrieved context.

% -------------------------------------------------------
\subsection{TounsiLM-8b Interface}
\label{subsec:llm}

The \textbf{TounsiLM-8b} model (\texttt{alabenayed/TounsiLM-8b}) is an 8-billion-parameter
LLM fine-tuned for understanding and generating Tunisian dialect using the
Llama-2 instruction format. The configurable generation parameters are summarized
in Table~\ref{tab:gen_params}.

\begin{table}[htbp]
\centering
\begin{tabular}{@{} l c l @{}}
\toprule
\textbf{Parameter}         & \textbf{Default} & \textbf{Role} \\
\midrule
\texttt{max\_new\_tokens}    & 512   & Maximum response length \\
\texttt{temperature}         & 0.7   & Creativity / diversity of generated text \\
\texttt{top\_p}              & 0.9   & Nucleus sampling \\
\texttt{top\_k}              & 50    & Top-$k$ sampling \\
\texttt{repetition\_penalty} & 1.1   & Repetition penalty \\
\bottomrule
\end{tabular}
\caption{TounsiLM-8b generation parameters.}
\label{tab:gen_params}
\end{table}

% -------------------------------------------------------
\subsection{Summary and Future Work}
\label{subsec:summary}

The developed RAG module brings several contributions to the TounsiLM system:

\begin{itemize}
    \item \textbf{Factual grounding}: responses are enriched with definitions,
    examples, and usage contexts drawn directly from the knowledge base, reducing
    hallucinations on dialectal vocabulary.

    \item \textbf{Arabizi robustness}: the query rewriter normalizes digit-based
    Arabizi spellings and generates multiple query variants, significantly improving
    recall for non-standard user inputs.

    \item \textbf{Hybrid precision}: combining BM25 and semantic search via RRF
    captures both exact keyword matches and semantic similarity, outperforming
    either method alone.

    \item \textbf{Automatic routing}: type-keyword detection routes queries to the
    appropriate KB category without user intervention, improving precision on
    domain-specific questions.

    \item \textbf{Extensibility}: adding a new data type requires only creating a
    Pydantic schema and a single line in the \texttt{SCHEMA\_MAP} dispatch table.

    \item \textbf{Transparency}: confidence level and sources are systematically
    exposed to the user, facilitating qualitative evaluation of responses.
\end{itemize}

Among the planned improvements are a \textbf{re-ranking} step using a cross-encoder
to further refine the top-$k$ results before generation, and a
\textbf{history-aware} conversational wrapper to maintain context across multiple
turns in interactive mode.

%  Required packages in the preamble:
%    \usepackage[utf8]{inputenc}
%    \usepackage[T1]{fontenc}
%    \usepackage[english]{babel}
%    \usepackage{graphicx}
%    \usepackage{booktabs}
%    \usepackage{tabularx}
%    \usepackage{listings}
%    \usepackage{xcolor}
%    \usepackage{tikz}
%    \usetikzlibrary{shapes.geometric, arrows.meta, positioning, fit}
%    \usepackage{amsmath}
%    \usepackage{hyperref}
% ============================================================

\section{Retrieval-Augmented Generation (RAG) Module}
\label{sec:rag}

\subsection{Introduction and Motivation}

Large language models (LLMs) suffer from a fundamental limitation: their knowledge is
frozen at training time and does not cover niche domains or low-resource languages.
Tunisian Arabic (\textit{Tounsi}) is precisely such a domain — poorly documented,
rich in loanwords from French, Classical Arabic and Berber, and subject to strong
regional and generational variation.

To address this limitation, we integrated a \textbf{Retrieval-Augmented Generation (RAG)}
module into the \textbf{TounsiLM-8b} model. The core idea is straightforward: before
generating a response, the system dynamically \textit{retrieves} the most relevant
entries from a specialized knowledge base and injects them into the LLM's context.
This grounds the generation in verified facts without requiring any retraining of the model.

% -------------------------------------------------------
\subsection{Overall Architecture}
\label{subsec:architecture}

The RAG pipeline is composed of six main stages, illustrated in
Figure~\ref{fig:rag_pipeline}.

\begin{figure}[htbp]
\centering
\begin{tikzpicture}[
    box/.style={rectangle, rounded corners=4pt, draw=black!70, fill=blue!8,
                minimum width=3.0cm, minimum height=0.85cm, align=center, font=\small},
    newbox/.style={rectangle, rounded corners=4pt, draw=black!70, fill=purple!10,
                minimum width=3.0cm, minimum height=0.85cm, align=center, font=\small},
    store/.style={cylinder, shape border rotate=90, draw=black!70, fill=orange!12,
                  minimum width=2.2cm, minimum height=1.3cm, align=center, font=\small},
    arrow/.style={-Stealth, thick},
    dasharrow/.style={-Stealth, thick, dashed},
    lbl/.style={font=\footnotesize\itshape, text=gray}
]

%--- Top row: inference flow (left -> right) ---%
\node[box]    (query)    at ( 0.0,  0.0) {User query\\{\small Arabic/Arabizi/French}};
\node[newbox] (rewriter) at ( 4.0,  0.0) {Query Rewriter\\+ Type Router};
\node[newbox] (hybrid)   at ( 8.0,  0.0) {Hybrid Retrieval\\BM25 + Semantic};
\node[store]  (chroma)   at (11.5,  0.0) {ChromaDB\\(vectors)};

%--- Bottom row: inference flow (right -> left) ---%
\node[box]    (merge)    at (11.5, -3.2) {Multi-query merge\\(best score)};
\node[box]    (context)  at ( 8.0, -3.2) {Context\\formatting};
\node[box]    (llm)      at ( 4.0, -3.2) {TounsiLM-8b\\(generation)};
\node[box]    (answer)   at ( 0.0, -3.2) {Answer +\\confidence level};

%--- Indexing column above ChromaDB (dashed) ---%
\node[box, fill=green!8] (indexer) at (11.5,  1.8) {Indexing\\(\texttt{passage:})};
\node[box, fill=green!8] (kb)      at (11.5,  3.3) {Knowledge base\\(JSON --- 1{,}647 entries)};

%--- Inference arrows ---%
\draw[arrow] (query)         -- (rewriter);
\draw[arrow] (rewriter)      -- node[above, lbl]{variants} (hybrid);
\draw[arrow] (hybrid)        -- (chroma);
\draw[arrow] (chroma.south)  -- (merge.north);
\draw[arrow] (merge)         -- (context);
\draw[arrow] (context)       -- (llm);
\draw[arrow] (llm)           -- (answer);
% Original question fed directly to LLM (left-side bypass)
\draw[arrow] (query.west) -- ++(-0.45,0) |- (llm.west);

%--- Indexing arrows ---%
\draw[dasharrow] (kb)      -- (indexer);
\draw[dasharrow] (indexer) -- node[right, lbl]{index} (chroma.north);

\end{tikzpicture}
\caption{RAG pipeline architecture. Purple boxes are newly introduced components.
Dashed path: one-time indexing phase. Solid path: inference flow.
The left-side arrow shows the original query fed directly to TounsiLM-8b alongside the retrieved context.}
\label{fig:rag_pipeline}
\end{figure}

% -------------------------------------------------------
\subsection{Knowledge Base}
\label{subsec:kb}

The knowledge base consists of JSON files organized by semantic category.
Every entry is validated against a Pydantic schema before indexing, ensuring
data integrity throughout the pipeline.

\subsubsection{Entry Types}

Table~\ref{tab:types} summarizes the eleven supported entry types, their associated
Pydantic schema, and the type-specific fields they introduce beyond the base schema.

\begin{table}[htbp]
\centering
\renewcommand{\arraystretch}{1.25}
\begin{tabularx}{\textwidth}{@{} l l X @{}}
\toprule
\textbf{Type} & \textbf{Schema} & \textbf{Type-specific fields} \\
\midrule
\texttt{expression}    & \texttt{ExpressionEntry}   & \texttt{origin}, \texttt{severity}, \texttt{gender\_sensitive} \\
\texttt{number\_slang} & \texttt{NumberSlangEntry}  & \texttt{msa\_equivalent} \\
\texttt{proverb}       & \texttt{ProverbEntry}      & \texttt{literal\_meaning}, \texttt{real\_meaning}, \texttt{when\_used}, \texttt{msa\_equivalent} \\
\texttt{food}          & \texttt{FoodEntry}         & \texttt{description}, \texttt{regional\_variation}, \texttt{when\_eaten}, \texttt{similar\_to} \\
\texttt{ingredient}    & \texttt{FoodEntry}         & (same as \texttt{food}) \\
\texttt{ritual}        & \texttt{RitualEntry}       & \texttt{occasion}, \texttt{expected\_response}, \texttt{tone} \\
\texttt{code\_switch}  & \texttt{CodeSwitchEntry}   & \texttt{origin\_language}, \texttt{origin\_word}, \texttt{tunisian\_equivalent}, \texttt{domain} \\
\texttt{tv\_series}    & \texttt{MediaEntry}        & \texttt{channel\_platform}, \texttt{era}, \texttt{cultural\_significance}, \texttt{common\_references} \\
\texttt{movie}         & \texttt{MediaEntry}        & (same as \texttt{tv\_series}) \\
\texttt{film}          & \texttt{MediaEntry}        & (same as \texttt{tv\_series}) \\
\texttt{color}         & \texttt{ColorEntry}        & \texttt{color\_family}, \texttt{cultural\_significance} \\
\bottomrule
\end{tabularx}
\caption{Knowledge base entry types and their Pydantic schemas.}
\label{tab:types}
\end{table}

\subsubsection{Base Schema}

All types inherit from \texttt{BaseEntry}, which enforces the following common fields:
\texttt{id}, \texttt{type}, \texttt{term\_arabic}, \texttt{term\_arabizi}, \texttt{meaning},
\texttt{meaning\_fr} (optional), \texttt{example}, \texttt{usage\_context},
\texttt{region}, \texttt{register}, \texttt{generation}, \texttt{scripts},
\texttt{source}, and \texttt{last\_updated}.

\subsubsection{Data Volume}

The current version of the knowledge base contains \textbf{1,647 valid entries}
across 9 JSON files, covering idiomatic expressions, proverbs, gastronomy, social
rituals, French--Arabic code-switching, media culture, and colors in the Tunisian dialect.

% -------------------------------------------------------
\subsection{Indexing Pipeline}
\label{subsec:indexing}

Indexing is performed in two steps: embedding text construction and insertion into
the vector store.

\subsubsection{Embedding Text Construction}

The \texttt{build\_embed\_text} module constructs a text for each entry by concatenating
its most informative fields. The selected fields vary by type: for a proverb,
\texttt{literal\_meaning} and \texttt{real\_meaning} are appended; for a
\texttt{code\_switch} entry, \texttt{origin\_language} and \texttt{tunisian\_equivalent}
are added.

Following the recommendations of the \texttt{multilingual-e5} model, every indexed
text is prefixed with \texttt{"passage: "}:

\begin{equation}
    t_{\text{index}} = \texttt{"passage: "} \;\|\; \bigoplus_{i} \text{field}_i
    \label{eq:passage_prefix}
\end{equation}

\noindent where $\|$ denotes concatenation and $\bigoplus$ the union of relevant fields.
This asymmetric prefix (\texttt{"passage:"} at indexing time, \texttt{"query:"} at
retrieval time) is essential for correctly exploiting the model's embedding space.

\subsubsection{Vector Store}

Vectors are persisted in \textbf{ChromaDB} using cosine similarity as the distance
metric. Indexing is incremental: identifiers already present in the collection are
skipped during updates. A deduplication mechanism using a \texttt{seen\_ids} set
prevents ID collisions across distinct JSON files.

% -------------------------------------------------------
\subsection{Query Rewriting and Routing}
\label{subsec:rewriting}

Arabizi --- the Latin-script transcription of Arabic widely used in online communication ---
has no standardized spelling. The same word can appear as \textit{barcha},
\textit{barca}, or \textit{bar5a}; \textit{7amra} and \textit{hamra} denote the same
color. A single query formulation therefore risks missing relevant entries that were
indexed under a different spelling.

\subsubsection{Query Rewriting}

The \texttt{QueryRewriter} module addresses this by generating up to three variants
of each query before retrieval:

\begin{enumerate}
    \item \textbf{Original query} --- preserved as-is.
    \item \textbf{Arabizi-normalized variant} --- common digit-to-letter substitutions
    are applied (Table~\ref{tab:arabizi}), converting e.g.\ \textit{7amra} to
    \textit{hamra} and \textit{5adra} to \textit{khadra}.
    \item \textbf{Cleaned variant} --- punctuation is stripped to handle
    question marks, commas, and Arabic \textit{tatweel}.
\end{enumerate}

\begin{table}[htbp]
\centering
\begin{tabular}{@{} c c l @{}}
\toprule
\textbf{Digit} & \textbf{Replacement} & \textbf{Arabic phoneme} \\
\midrule
3 & a        & ayn \\
7 & h        & ha \\
5 & kh       & kha \\
9 & q        & qaf \\
8 & gh       & ghain \\
2 & (dropped) & hamza \\
\bottomrule
\end{tabular}
\caption{Arabizi digit-to-letter normalization rules applied by the query rewriter.}
\label{tab:arabizi}
\end{table}

Each variant is retrieved independently. Results are then merged by keeping the
highest score for each document across all variants:

\begin{equation}
    s^*(d) = \max_{v \in \mathcal{V}} s(d, v)
    \label{eq:multiquery_merge}
\end{equation}

\noindent where $\mathcal{V}$ is the set of query variants and $s(d,v)$ is the
RRF-fused score of document $d$ for variant $v$.

\subsubsection{Query Routing}

The same \texttt{QueryRewriter} module inspects the query for type-specific keywords
and automatically sets the \texttt{entry\_type} filter before retrieval. For example,
a query containing \textit{harissa}, \textit{food}, or the Arabic term for cooking
routes to \texttt{entry\_type="food"}, restricting the search space to the relevant
category. This improves both precision and speed without any user intervention.

% -------------------------------------------------------
\subsection{Hybrid Retrieval}
\label{subsec:retrieval}

Retrieval combines two complementary methods: dense semantic search and sparse
keyword search (BM25). Each handles a different failure mode of the other.

\subsubsection{Semantic Retrieval}

At inference time, the query variant is prefixed with \texttt{"query: "} and encoded
by \texttt{multilingual-e5-base}:

\begin{equation}
    \mathbf{q} = \text{Encode}\!\left(\texttt{"query: "} \| q\right)
    \label{eq:query_encode}
\end{equation}

ChromaDB computes cosine similarity between $\mathbf{q}$ and all indexed passage
vectors, returning $3k$ candidates to leave room for RRF re-ordering. The relevance
score is:

\begin{equation}
    s_{\text{sem}}(d) = 1 - d_{\cos}(\mathbf{q},\, \mathbf{p}_d)
    \label{eq:sem_score}
\end{equation}

\subsubsection{BM25 Retrieval}

A \textbf{BM25Okapi} index is built at startup over all 1,647 documents (stripped of
the \texttt{"passage: "} prefix). At query time, BM25 scores are normalized to $[0,1]$:

\begin{equation}
    s_{\text{BM25}}(d) = \frac{\text{BM25}(q, d)}{\max_{d'} \text{BM25}(q, d')}
    \label{eq:bm25_score}
\end{equation}

BM25 excels at exact term matches --- when a user types a precise Arabizi word or a
French loanword that dense embeddings might misplace in the embedding space.

\subsubsection{Reciprocal Rank Fusion}

The two ranked lists are merged using \textbf{Reciprocal Rank Fusion (RRF)} with
per-source weights:

\begin{equation}
    \text{RRF}(d) = \frac{w_{\text{sem}}}{k_0 + r_{\text{sem}}(d)}
                  + \frac{w_{\text{BM25}}}{k_0 + r_{\text{BM25}}(d)}
    \label{eq:rrf}
\end{equation}

\noindent where $r_i(d)$ is the rank of document $d$ in ranking $i$ (1-indexed),
$k_0 = 60$ is the RRF smoothing constant, $w_{\text{sem}} = 0.6$ and
$w_{\text{BM25}} = 0.4$. The final score is normalized to $[0,1]$ and the top-$k$
documents are returned.

% -------------------------------------------------------
\subsection{Confidence Scoring}
\label{subsec:confidence}

After retrieval, an overall confidence score is computed from the mean of the
RRF-normalized scores:

\begin{equation}
    \bar{s} = \frac{1}{k} \sum_{i=1}^{k} s_i
    \label{eq:mean_score}
\end{equation}

The confidence level is then assigned according to three tiers:

\begin{table}[htbp]
\centering
\begin{tabular}{@{} c l l @{}}
\toprule
\textbf{Indicator} & \textbf{Level} & \textbf{Condition} \\
\midrule
\textcolor{green!60!black}{$\bullet$} & High   & $\bar{s} \geq 0.75$ \\
\textcolor{orange}{$\bullet$}         & Medium & $0.50 \leq \bar{s} < 0.75$ \\
\textcolor{red}{$\bullet$}            & Low    & $\bar{s} < 0.50$ \\
$\bullet$                             & None   & no results returned \\
\bottomrule
\end{tabular}
\caption{Confidence levels based on the mean RRF score $\bar{s}$.}
\label{tab:confidence}
\end{table}

% -------------------------------------------------------
\subsection{Context Formatting}
\label{subsec:context}

The retrieved entries are formatted into a structured text block injected into the
LLM prompt. Each line follows the format:

\begin{quote}
\texttt{[type | score: \textit{s}] arabic\_term (arabizi) : meaning | \textnormal{مثال:} example | \textnormal{الاستخدام:} context}
\end{quote}

To avoid exceeding the model's context window (4,096 tokens), truncation is performed
at the token level using the TounsiLM-8b tokenizer, with a default limit of
\textbf{1,500 tokens} for the retrieved context.

% -------------------------------------------------------
\subsection{TounsiLM-8b Interface}
\label{subsec:llm}

The \textbf{TounsiLM-8b} model (\texttt{alabenayed/TounsiLM-8b}) is an 8-billion-parameter
LLM fine-tuned for understanding and generating Tunisian dialect using the
Llama-2 instruction format. The configurable generation parameters are summarized
in Table~\ref{tab:gen_params}.

\begin{table}[htbp]
\centering
\begin{tabular}{@{} l c l @{}}
\toprule
\textbf{Parameter}         & \textbf{Default} & \textbf{Role} \\
\midrule
\texttt{max\_new\_tokens}    & 512   & Maximum response length \\
\texttt{temperature}         & 0.7   & Creativity / diversity of generated text \\
\texttt{top\_p}              & 0.9   & Nucleus sampling \\
\texttt{top\_k}              & 50    & Top-$k$ sampling \\
\texttt{repetition\_penalty} & 1.1   & Repetition penalty \\
\bottomrule
\end{tabular}
\caption{TounsiLM-8b generation parameters.}
\label{tab:gen_params}
\end{table}

% -------------------------------------------------------
\subsection{Summary and Future Work}
\label{subsec:summary}

The developed RAG module brings several contributions to the TounsiLM system:

\begin{itemize}
    \item \textbf{Factual grounding}: responses are enriched with definitions,
    examples, and usage contexts drawn directly from the knowledge base, reducing
    hallucinations on dialectal vocabulary.

    \item \textbf{Arabizi robustness}: the query rewriter normalizes digit-based
    Arabizi spellings and generates multiple query variants, significantly improving
    recall for non-standard user inputs.

    \item \textbf{Hybrid precision}: combining BM25 and semantic search via RRF
    captures both exact keyword matches and semantic similarity, outperforming
    either method alone.

    \item \textbf{Automatic routing}: type-keyword detection routes queries to the
    appropriate KB category without user intervention, improving precision on
    domain-specific questions.

    \item \textbf{Extensibility}: adding a new data type requires only creating a
    Pydantic schema and a single line in the \texttt{SCHEMA\_MAP} dispatch table.

    \item \textbf{Transparency}: confidence level and sources are systematically
    exposed to the user, facilitating qualitative evaluation of responses.
\end{itemize}

Among the planned improvements are a \textbf{re-ranking} step using a cross-encoder
to further refine the top-$k$ results before generation, and a
\textbf{history-aware} conversational wrapper to maintain context across multiple
turns in interactive mode.
